% !TeX program = xelatex
\documentclass[11pt,a4paper]{article}

\usepackage[margin=22mm]{geometry}
\usepackage{amsmath}
\usepackage{array}
\usepackage{booktabs}
\usepackage{graphicx}
\usepackage{hyperref}
\usepackage{iftex}
\usepackage{tabularx}
\usepackage{xcolor}

\ifPDFTeX
  \usepackage[utf8]{inputenc}
  \usepackage{CJKutf8}
  \newcommand{\startjapanese}{\begin{CJK*}{UTF8}{min}}
  \newcommand{\stopjapanese}{\end{CJK*}}
\else
  \usepackage{fontspec}
  \IfFontExistsTF{Noto Serif CJK JP}{
    \setmainfont{Noto Serif CJK JP}
  }{
    \setmainfont{Hiragino Mincho ProN}
  }
  \IfFontExistsTF{Noto Sans CJK JP}{
    \setsansfont{Noto Sans CJK JP}
  }{
    \setsansfont{Hiragino Kaku Gothic ProN}
  }
  \ifXeTeX
    \XeTeXlinebreaklocale "ja"
    \XeTeXlinebreakskip=0pt plus 1pt
  \fi
  \newcommand{\startjapanese}{}
  \newcommand{\stopjapanese}{}
\fi

\definecolor{oublue}{HTML}{0034B8}
\definecolor{darknavy}{HTML}{070D2F}
\definecolor{softblue}{HTML}{F2F6FF}
\definecolor{warningred}{HTML}{C81E1E}
\definecolor{mutedgray}{HTML}{596173}

\ifPDFTeX
  \hypersetup{
    colorlinks=true,
    linkcolor=oublue,
    urlcolor=oublue,
    pdfauthor={Osaka University, Group 5},
    pdftitle={Development of OU Textbook}
  }
\else
  \hypersetup{
    colorlinks=true,
    linkcolor=oublue,
    urlcolor=oublue,
    pdfauthor={大阪大学 基礎工学部 情報科学科 5班},
    pdftitle={大阪大学生限定 参考書リユース取引アプリ OU Textbook の開発}
  }
\fi

\setlength{\parindent}{0pt}
\setlength{\parskip}{0.65em}
\renewcommand{\arraystretch}{1.25}

\newcommand{\oustrong}[1]{\textbf{\color{darknavy}#1}}
\newcommand{\flowtitle}[1]{\textbf{\color{oublue}#1}}

\title{
  \vspace{-1.5cm}
  {\color{oublue}\LARGE\bfseries 大阪大学生限定}\\[0.35em]
  {\Huge\bfseries 参考書リユース取引アプリ}\\[0.25em]
  {\Huge\bfseries 「OU Textbook」の開発}
}

\author{
  大阪大学 基礎工学部 情報科学科 \quad 5班\\[0.8em]
  \small
  \begin{tabular}{lll}
    \toprule
    学籍番号 & 氏名 & 電子メールアドレス\\
    \midrule
    09B25048 & 杉山 優弥 & \texttt{u966882c@ecs.osaka-u.ac.jp}\\
    09B25064 & 沼田 智也 & \texttt{u493047f@ecs.osaka-u.ac.jp}\\
    09B25073 & 福田 侑飛 & \texttt{u649863a@ecs.osaka-u.ac.jp}\\
    09B25080 & 正木 光一郎 & \texttt{u756057d@ecs.osaka-u.ac.jp}\\
    09B25081 & 益澤 悠 & \texttt{u315289b@ecs.osaka-u.ac.jp}\\
    09B25099 & 山下 初一音 & \texttt{u267690f@ecs.osaka-u.ac.jp}\\
    \bottomrule
  \end{tabular}
}

\date{\today}

\begin{document}
\startjapanese

\maketitle
\thispagestyle{empty}

\section{目的}

大阪大学の学生が、授業で使い終わった教科書や参考書を、学内で手軽に売買・譲渡できるスマホ・Webアプリを開発することを目的とした。

商品の検索、購入相談、価格交渉、受け渡し日時・場所の決定、取引完了までを一つのアプリ内で行えるようにし、学生の教科書購入費の負担と、使われなくなった教科書の廃棄を減らすことを目指した。

\section{動機}

大学では、授業が終わるたびに使わなくなった教科書が増え、整理や処分に困る学生がいる。教科書は高価なものも多いため、処分するのではなく、できれば売りたい人や、必要な学生へ譲りたい人もいる。

また、専門とは異なる授業のために購入した高価な教科書が、授業でほとんど使われず、新品に近い状態のまま残ってしまう場合もある。そのような教科書を捨ててしまうのは非常にもったいないと考えた。

一方、一般的なフリマサービスでは、送料や手数料がかかるほか、発送のために教科書を梱包し、配送業者へ持参する必要がある。

そこで、大阪大学の学生同士であれば、普段利用するキャンパス内で直接受け渡しができ、送料をかけずに、教科書を必要とする学生へ引き継げると考えた。

\section{アプローチ・工夫}

\begin{itemize}
  \item 大阪大学のメールアドレスを利用し、利用者を大学関係者に限定した。登録時に入力されたメールアドレスへ認証メールを送信し、メール内のリンクから本人確認を完了するメール認証機能を実装した。
  \item 利用開始時の個人情報の取扱いと取引上のルールを明確にするため、新規登録画面から利用規約とプライバシーポリシーを確認できるようにした。両方への同意を必須とし、未同意の場合は認証メールの送信および登録処理へ進めないよう、画面とサーバーの両方で制御した。
  \item 書名、著者、教材区分、受け渡し希望キャンパス、商品の状態、価格帯による検索・絞り込み機能を実装した。
  \item 出品、画像登録、お気に入り、購入相談、取引メッセージを一つのサービスに統合した。
  \item 出品者が価格を提示し、購入希望者が同意して初めて取引相手が確定する仕組みを実装した。
  \item 価格だけでなく、受け渡し日時と場所も双方の同意によって決定するようにした。
  \item 一人との取引が成立した後は、他の利用者が同じ商品へ同意できないように制御した。
  \item 出品中、取引中、受け渡し待ち、売却済み、取り下げ済みなど、商品の状態を段階的に管理した。
  \item 安心して取引相手を選べるように、受け渡し完了後に出品者と購入者が互いを評価する相互評価機能を実装した。片方の評価だけで点数が変わらないようにし、双方が評価を送信した時点で信用点数へ反映する。
  \item 商品画像をクラウドストレージへ保存し、アプリの再起動や再デプロイ後も公開環境で継続して表示できる構成を採用した。
  \item \oustrong{スマートフォン・タブレット・PCへのレスポンシブ対応}を行った。スマートフォンでの操作性を維持しながら、タブレットでは一覧を複数列で表示し、PCでは下部ナビゲーションを左サイドナビゲーションへ切り替えることで、画面幅を有効に使い、端末に応じて見やすく操作しやすい画面構成にした。
\end{itemize}

\subsection{信用点数制度}

信用点数は100点から始まり、取引時の評価やキャンセル状況に応じて増減する。

\begin{table}[htbp]
  \centering
  \caption{行動に応じた信用点数の変化}
  \begin{tabularx}{0.86\textwidth}{Xr}
    \toprule
    行動 & 点数の変化\\
    \midrule
    取引相手から良い評価を受ける & $+10$点\\
    取引相手から悪い評価を受ける & $-10$点\\
    受け渡し予定の確定後に自己都合でキャンセルする & $-10$点\\
    連絡なく受け渡しに来なかったと報告される & $-30$点\\
    \bottomrule
  \end{tabularx}
\end{table}

\begin{table}[htbp]
  \centering
  \caption{信用点数と信用ランク}
  \begin{tabular}{cl}
    \toprule
    信用点数 & ランク\\
    \midrule
    150点以上 & エキスパート\\
    120点以上149点以下 & トラスト\\
    41点以上119点以下 & レギュラー\\
    40点以下 & \textcolor{warningred}{\textbf{要注意}}\\
    \bottomrule
  \end{tabular}
\end{table}

例えば、120点の利用者が悪い評価を受けて110点になると「トラスト」から「レギュラー」へ、150点の利用者が悪い評価を受けて140点になると「エキスパート」から「トラスト」へ下がる。また、信用点数が40点以下になると、取引時に注意が必要な利用者として「要注意」と表示する。

\section{難しかった点・頑張った点}

\subsection{大阪大学生に限定したメール認証}

Supabaseのメール認証を活用し、大阪大学のメールアドレスを持つ本人であることを確認できる仕組みを実装した。登録時に認証メールを送信し、メール内のリンクを開いた利用者だけが登録を完了できるようにした。また、認証メールを送信する前に利用規約とプライバシーポリシーの確認・同意を必須とし、未同意の登録をサーバー側でも拒否することで、単なる画面上のチェックだけにならないようにした。

\subsection{取引フローと商品状態の管理}

最も難しかったのは、取引の流れと商品の状態を正しく管理することである。複数の購入希望者が同じ商品について相談できる一方で、最終的に取引できるのは一人だけである。そのため、価格への同意、受け渡し条件の決定、キャンセル、取引完了といった操作が競合しても、商品の状態に矛盾が起きないように設計した。

実装中には、取り下げた商品が「取り下げ一覧」に表示されず「出品一覧」に残る、取引中になった商品が「取引中／売却済み一覧」に表示されないなど、内部の状態と画面表示が一致しない問題が発生した。商品の状態を更新する処理と、各一覧で商品を抽出する条件の両方を見直し、操作後に正しい一覧へ表示されるよう改善した。

また、画像アップロード、通知、メッセージ、マイページの表示を連動させ、利用者が現在の取引状況を迷わず確認できるように改善を重ねた。

\section{結果}

利用登録から参考書の受け渡し、取引相手の評価までを、一つのWebアプリ内で完結できるようにした。また、\oustrong{レスポンシブ対応}により、スマートフォン、タブレット、PCの画面幅に応じてレイアウトが切り替わり、どの端末からでも主要な機能を利用できるようにした。

\begin{enumerate}
  \item \flowtitle{大阪大学メールによる利用登録}\\
  利用者は名前と大阪大学のメールアドレスを入力し、利用規約とプライバシーポリシーを確認して同意した上で登録する。両方に同意しない場合は登録処理へ進めない。入力されたアドレスへ認証メールを送り、メール内のリンクを開いて本人確認を完了した学生だけが利用できる。登録後は、学部・学科・学年をプロフィールに設定できる。

  \item \flowtitle{参考書の出品}\\
  出品者は、書名、著者、価格、教材区分、商品の状態、受け渡しを希望するキャンパス、説明、商品画像を登録する。画像はクラウドストレージへ保存されるため、アプリを更新・再起動しても継続して表示される。

  \item \flowtitle{参考書の検索と絞り込み}\\
  購入希望者は、書名や著者などのキーワードで参考書を検索できる。さらに、教材区分、受け渡し希望キャンパス、商品の状態、価格帯で絞り込み、新着順、価格順、お気に入り数順に並べ替えられる。

  \item \flowtitle{お気に入りと出品者情報の確認}\\
  気になる参考書をお気に入りへ登録し、後からマイページで確認できる。商品詳細では出品者名を選択し、信用点数、信用ランク、出品数、取引完了数を確認できる。

  \item \flowtitle{購入相談と取引メッセージ}\\
  購入希望者は商品詳細から出品者とのチャットを開始し、商品の状態や受け渡しについて相談できる。一つの商品に複数の購入希望者が相談できるが、チャットは相手ごとに分けて管理される。

  \item \flowtitle{取引価格の提示と合意}\\
  出品者はチャット内で取引価格を提示する。購入希望者は、提示価格に「同意する」または「同意しない」を選択できる。同意された時点で取引相手を一人に確定し、他の購入希望者が同じ商品へ同意できないように制御する。

  \item \flowtitle{受け渡し日時・場所の提示と合意}\\
  取引相手の確定後、出品者が受け渡し日時と場所を提示する。購入者は、その条件にも「同意する」または「同意しない」を選択できる。双方の合意後は、受信箱とマイページに確定した日時・場所を目立つ形で表示する。

  \item \flowtitle{学内での受け渡し}\\
  確定した日時・場所で参考書を直接受け渡す。受け渡し予定時刻になると、出品者と購入者の双方が「受け渡し完了」を記録できるようになる。

  \item \flowtitle{双方による完了確認}\\
  片方が完了ボタンを押しただけでは取引を完了させず、もう一方の確認を待つ。双方が完了を確認すると、それぞれへ完了通知を送り、取引相手の評価画面へ進む。

  \item \flowtitle{相互評価と取引完了}\\
  出品者と購入者が互いに「良い評価」または「悪い評価」を送信する。双方の評価がそろった時点で信用点数へ反映し、商品を売却済みにして一連の取引を完了する。

  \item \flowtitle{取り下げ・キャンセルへの対応}\\
  出品者は取引開始前の商品を取り下げ、取り下げ一覧から確認できる。取引が途中でキャンセルされた場合は、商品を再び出品中へ戻す。また、自己都合のキャンセルや無断キャンセルは信用点数へ反映される。
\end{enumerate}

これらの主要機能に対する自動テストも作成し、登録、出品、検索、条件合意、受け渡し、評価までの動作を継続的に確認できるようにした。

\section{他にアピールしたいこと}

本アプリは、一般的なフリマサイトを小さく再現したものではなく、\oustrong{大阪大学の学生が学内で参考書を直接受け渡す状況}に合わせて設計している。

特に、価格・日時・場所を双方の同意で確定する仕組みと、複数の購入希望者がいる場合でも取引相手を一人に限定する仕組みが中心的な特徴である。

今後は、授業名やシラバスとの連携、履修科目に応じた参考書の推薦、取引実績を活用した信頼性向上などを検討している。

\stopjapanese
\end{document}
